% Self-contained source composed from pinned proofs and measured outputs.
\documentclass[12pt]{article}
\usepackage[margin=0.75in]{geometry}
\usepackage[T1]{fontenc}
\usepackage{amsmath,amssymb,amsthm,booktabs,array}
\usepackage[authoryear,round]{natbib}
\usepackage[hidelinks]{hyperref}
\usepackage{xurl,microtype,placeins}
\usepackage[mathlines]{lineno}
\AtBeginDocument{\linenumbers}
\setlength{\emergencystretch}{2em}
\renewcommand{\topfraction}{0.95}
\renewcommand{\bottomfraction}{0.9}
\renewcommand{\textfraction}{0.05}
\renewcommand{\floatpagefraction}{0.75}
\newtheorem{theorem}{Theorem}
\newtheorem{proposition}{Proposition}
\newtheorem{corollary}{Corollary}
\newcommand{\F}{\mathcal{F}}
\newcommand{\R}{\mathbb R}
\newcommand{\M}{\mathcal M}
\newcommand{\Oset}{\mathcal O}
\newcommand{\clip}{\operatorname{clip}}
\title{Risk and Adaptivity in Certifying Machine-Learning Model Comparisons}
\author{Jishan Kapoor\\Independent Researcher\\Toronto M4Y 1C9, Canada\\
Corresponding author: \href{mailto:kapoorjishan19@gmail.com}{kapoorjishan19@gmail.com}}
\date{}
\begin{document}
\maketitle
\begin{abstract}
\textbf{Background.} Auditing a machine-learning model can require retrieving
missing verification targets. The auditor must choose between a full export,
a fixed batch of measurements and queries that depend on earlier answers.
These choices can have different guarantees even with the same reading budget.

\textbf{Methods.} We study the sharp interval containing a model-loss contrast
over all records consistent with the available information. The contrast is
affine, the uncertainty set compact and convex, and the query dictionary finite
and linear. We classify terminal-width penalties and give exact two-query
algorithms with independently checkable certificates.

\textbf{Results.} Randomized batch and adaptive minimax risks agree for every
such problem if and only if the nondecreasing penalty is concave. Constructed
families have identical optimal expected widths but sharply different precision
guarantees. In 290 of 4,914 frozen forecast comparisons, two adaptive counter
readings guarantee halving the initial width, whereas optimal randomized-batch
worst-case success is 50\%. An external source reproduces this gap in 29 of
1,674 comparisons. Overall recorded gains over the best deterministic
worst-width batch are 0.407 and 0.119 percentage points. Under illustrative
1\% internal-reading errors with exact endpoints, all 29 external cases remain
certifiable using two readings in 22 cases and three in seven. Capacity bounds
often remove feedback necessity, and measured public-API retrieval favors bulk.

\textbf{Conclusions.} The method identifies which acquisition protocol meets
a declared assurance requirement. An executable selector chooses among supplied certified protocols using
declared measurement assumptions and resource limits.
\end{abstract}
\noindent\textbf{Keywords:} artificial intelligence; machine learning;
adaptive information; minimax risk; optimal recovery; robust optimization.

\section{Introduction}
A model audit may need additional measurements before it can justify a
comparison. The auditor can retrieve the full record, choose a batch of
readings in advance, or choose later readings after seeing earlier answers.
We call the last option adaptive querying. A sharp interval is the smallest
interval containing the loss difference for every record compatible with
the available information. This paper asks when feedback improves the
guarantee that this interval will be sufficiently narrow. A narrow interval
can still contain zero, so precision and identifying the better model are
separate goals.

We study affine scalar contrasts on nonempty compact convex uncertainty
sets, with queries selected from a fixed finite dictionary of linear
measurements. This formulation includes paired squared-loss comparisons,
overlapping linear constraints, valid coordinate bounds, signed targets,
and fixed query costs. Batch and adaptive policies have the same hard
budget. The unknown record is fixed before the algorithm's random seed is
drawn; probabilities describe acquisition randomization, not a distribution
assumed for the missing truth. Fixed bounded measurement errors can be
included by enlarging the hidden record when the resulting uncertainty set
is compact and convex.

The computer-science question is how an acquisition algorithm should use
feedback when its output must be valid for every compatible record. We ask
which risk objectives permit a feedback advantage, how to compute an
instance-specific guarantee, and which permitted retrieval protocol should
implement it. The fitted forecasting models supply fixed prediction vectors;
the proposed algorithm governs their verification data, rather than learning
a new predictor.

The contribution is a classification of when feedback can change minimax
certificate risk, sharp realizations using admissible information queries,
and an exact two-query decision rule with checkable precision guarantees.
Concavity exactly characterizes the nondecreasing width penalties for which
adaptation never improves randomized minimax risk. A coordinate-query
construction proves the converse, and a family separates equal mean-width
optima from success probabilities $1/K$ and one. The classification uses
established minimax reasoning and Jensen's inequality. A dimension-dependent
precision consequence adapts established sparse-policy bounds
\citep{kurtz2026}; we do not claim a new general sparsification theorem.

Cumulative meter records test whether this distinction affects actual
forecast comparisons. Before optional answers are revealed, the two-query
rule identifies cases where a fixed batch suffices, feedback changes
feasibility, or two queries do not provide a deterministic guarantee.
The original and external studies contain 290 of 4,914 and 29 of 1,674 cases,
respectively, with an exact adaptive-versus-batch precision-success gap of
one versus one half. These cases span three meters and six model pairs in
each source. The large worst-case assurance gap is distinct from the small
overall recorded gains of 0.407 and 0.119 percentage points.

The acquisition interface determines which guarantee is useful. Under the
declared 1\% internal-reading error allowance and exact endpoints, all 29
selected external cases remain certifiable: 22 use two readings and seven
use three. Added capacity information can remove the need for feedback.
On the measured public API, bulk retrieval is preferable. Our implementation
therefore chooses among supplied certificates using the permitted
measurements, reading and round limits, and cost priority. It reports when
those certificates cannot meet the declared requirements.

The paper develops the classification and sharp examples, then computable
two-query guarantees and their cumulative-meter specialization. The
evaluation reports fixed protocols, recorded replay, uncertainty sensitivity
and actual retrieval measurements. Online Resource~1 provides extended
proofs, the separate compact coordinate-query solver, historical controls,
complete case ledgers and executable verification.

\subsection{Related work and the contribution boundary}
Partial identification studies conclusions that survive incomplete
observations \citep{manski,polo}. For model comparison, cancellation in
paired squared losses leaves an affine function of the unknown truth;
this is established loss structure rather than a new loss function
\citep{patton}. Active evaluation and model selection acquire labels for
estimating performance or choosing among predictors
\citep{sawade,hara,active,okanovic2025}. Our target is the entire sharp
interval of a fixed-record contrast after exact observations. No estimator
of the population risk, surrogate label distribution, or new forecasting
architecture is introduced.

% Krieg--Ullrich locators below refer explicitly to the author manuscript;
% journal pagination is not inferred from the manuscript page numbers.
The certificate width is the established local diameter of information
\citep[Section~4.1]{novak2008}.
\citet[Theorems~2--3 and~5]{osipenko} also studies when adaptation cannot
improve noisy scalar recovery. Those criteria concern deterministic
reconstruction error and the existence of an optimal nonadaptive linear
reconstruction; the present criterion instead classifies randomized
penalties of the complete terminal information-fiber width.
Classical deterministic results already establish strong limits on
adaptivity: on convex symmetric classes with linear information,
\citet[Theorem~4.4]{novak2008} construct nonadaptive information with
no greater global diameter and at most twice the reconstruction radius.
For scalar targets the radius comparison is exact
\citep[Theorem~4.5]{novak2008}. A modern formulation of the factor-two
bound appears in \citet[Proposition~4.1]{krieg2025}.
Nonsymmetric convex uncertainty itself is also established:
\citet{novak1995} compares adaptive and nonadaptive optimal-recovery
errors using relations between Gelfand and Bernstein widths. The
classical objective is reconstruction error, whereas our randomized
penalty applies to the full terminal compatible interval. Our classification allows nonsymmetric convex sets and concerns randomized
penalties of the terminal local diameter, including a converse and exact
precision and moment separations. Its mean-risk interchange applies
\citet[Theorem~3.4]{sion1958}; the concave-penalty upper bound uses Jensen's
inequality.

For fixed deterministic scalar information, the optimal worst-case
reconstruction error is half the largest fiber width. The distinction here
concerns randomized local-width penalties and the order of averaging and
worst-case evaluation. Specifically, \citet[Section~2]{krieg2025}
and \citet[Section~9.1, Eq.~(9.15), author manuscript v2]{krieg2026}
study $\inf_A\sup_{y\in F}\mathbb E_\omega
\lVert T(y)-A_\omega(y)\rVert$, the error of a selected reconstruction.
Here the random transcript instead incurs the full compatible-interval
width $W_{Q_\omega(y)}(y)$.
Writing its scalar local radius as $r_{Q_\omega(y)}(y)=W_{Q_\omega(y)}(y)/2$,
our mean risk is $2\sup_y\mathbb E_\omega r_{Q_\omega(y)}(y)$.
For a randomized batch, this generally differs from
$2\mathbb E_\omega\sup_y r_{Q_\omega}(y)$, twice the expected global radius
discussed in \citet[Eq.~(2.4), Remark~6.1 and Section~6.1, author manuscript v2]{krieg2026}.
The midpoint reconstruction has pointwise error at most the local radius,
but its error at the actual record need not equal that radius. Consequently,
neither replacing actual reconstruction error by interval width nor
interchanging the expectation and supremum preserves the optimization
problem. The present theorem concerns the terminal local-width risk.
The fixed measurement dictionary is substantive: if every linear
functional were directly queryable at the same cost, a scalar affine
target could be determined by querying its linear part once.


Information-dependent decisions are also established in adjustable robust
optimization \citep{wei2024,ddid}. The dimension-dependent precision
bound below adapts the sufficient-policy result of
\citet[Theorem~2 and Example~5]{kurtz2026}; its query interpretation and
sharp examples are made explicit. Weighted-oracle randomized optimization
and dynamic-program LP constructions are established
\citep{kawase2024,martin1990}. The supporting aggregate solver contributes
a structured hard-budget batch representation and dual-completion
certificate (Online Resource~1, S5/S7). Secondary optimization among tied
robust solutions follows an established concern \citep{iancu2014}.

Adaptive sampling of cumulative information has close precedents.
\citet[Section~4.1.1, Corollary~1 and Theorem~3]{gaillard2023}
study certified scalar integral estimation from sequential evaluations of
a monotone function. For monotone contrast coefficients, summation by parts
relates the register problem to weighted monotone integration. Our instance
test pairs an adaptive precision guarantee with an adverse distribution
limiting every randomized batch under the same information interface.

\citet[Theorem~1 and Proposition~1]{kuboyama2026} gives exact model-choice
certificates after adaptive binary labels and compares shared-label paired
bounds with separate intervals. That work also distinguishes hindsight
certificate size from acquisition requirements (Theorem~2 and Proposition~3).
Our application uses constrained continuous targets, while the central
classification concerns the penalty of terminal sharp width. These precise
distinctions, rather than certification or shared comparisons generally,
define the contribution relative to model-choice and information-discovery
frameworks.

\section{Methods}\label{sec:methods}
\subsection{Model Comparisons and Sharp Information}
Let $a,b\in\mathbb R^N$ be two fixed prediction vectors for the same
scheduled targets. Their paired mean-squared-error difference is
\begin{equation}
 D(y)=\frac1N\sum_{i=1}^N\{(a_i-y_i)^2-(b_i-y_i)^2\}
 =d+c^\top y,\qquad c_i=\frac{2(b_i-a_i)}N.
 \label{eq:linear}
\end{equation}
All scheduled targets remain in the denominator. Initially observed values,
documented aggregate measurements and other defensible constraints define
the feasible set $F$. Any permitted archive record that directly reconstructs
a target is incorporated before this set is declared. Sharpness is always
relative to that declared information set.

The attainable interval of $D$ on $F$ is an identification interval for
this record, not a sampling confidence interval. A negative upper endpoint
certifies that $a$ has smaller squared loss for every completion; a positive
lower endpoint certifies $b$. An interval containing zero cannot certify
a strict winner. Subsequent exact measurements restrict $F$ to a fiber and
therefore cannot widen its sharp interval. A prescribed width threshold
instead asks for a precision guarantee, regardless of the interval's location.

The mathematical examples below also describe paired squared-loss widths.
Given any contrast coefficient $c$, nonnegative predictions
$a_i=N\max(-c_i,0)/2$ and $b_i=N\max(c_i,0)/2$ realize that coefficient
in \eqref{eq:linear}. The induced constant changes the interval's position,
but not its width. Thus the constructions have a model-comparison
interpretation without being measurements of prediction performance.

\subsection{Information, terminal width, and risk}
\label{sec:risk_setup}
Let $F\subset\mathbb R^n$ be nonempty, compact, and convex, and let
$T(y)=c^\top y+d$ be the scalar quantity to be audited. A fixed finite
collection of exact linear measurements is available. Write $L_Q$ for
the measurement matrix of a selected set $Q$, and let $\mathcal Q$ be
the nonempty finite family of admissible final sets, available to both
policy classes. Admissibility is determined before observing answers; for example,
$\sum_{j\in Q}k_j\le B$ for nonnegative fixed costs $k_j$.
Coordinate observations correspond to rows of the identity matrix.
The results do not require nonnegative targets or disjoint aggregates.

For a feasible record $y$, define the information fiber and its sharp
target width by
\begin{equation}
 F_Q(y)=\{z\in F:L_Qz=L_Qy\},\qquad
 W_Q(y)=\max_{z\in F_Q(y)}T(z)-\min_{z\in F_Q(y)}T(z).
 \label{eq:risk_fiber}
\end{equation}
The extrema exist by compactness, and $0\le W_Q(y)\le D_T$, where
$D_T=\max_F T-\min_F T$. This is the local diameter of the
scalar information fiber. The optimized quantity is the full fiber's
diameter, rather than the actual error of a selected point estimate;
it is not a confidence interval. An additive constant in $T$ does
not affect $W_Q$. In particular, a paired squared-loss contrast is
affine in the unknown verification targets and fits this formulation.

A randomized batch policy draws $Q\sim\pi$ before seeing answers. An
adaptive policy selects successive measurements and a stopping decision
from the answers already observed and an independent internal random
seed. Every execution must terminate after finitely many queries at a
set in $\mathcal Q$. Query and stopping rules are jointly measurable in
the observed history and seed; observations are exact. For a finite-valued,
nondecreasing function $\phi:[0,\infty)\to\mathbb R$, define
\begin{equation}
 R_{\mathrm{batch}}^{\phi}
 =\inf_{\pi\in\Delta(\mathcal Q)}\sup_{y\in F}
   \mathbb E_{Q\sim\pi}\phi(W_Q(y)),\qquad
 R_{\mathrm{adaptive}}^{\phi}
 =\inf_{\mathcal P}\sup_{y\in F}
   \mathbb E\phi(W_{Q_{\mathcal P}(y)}(y)).
 \label{eq:risk_objectives}
\end{equation}
The hidden record is fixed before the random seed is drawn. Thus the
supremum is outside the expectation. When measurement error is represented
by additional hidden coordinates, one fixed error vector is part of that
record for the entire execution. An adversary choosing a new error after
each query is a different information model. These objectives concern terminal
width under a hard budget, rather than expected requests until a winner
is certified. The expectations are well defined: the widths are Borel
measurable by Proposition~\ref{prop:risk_concavity_transcript}, and a
finite-valued monotone $\phi$ is Borel measurable and bounded on $[0,D_T]$.

For a required precision $\tau\geq0$, let $\mathcal C$ be either the
randomized batch or randomized adaptive policy class, and write
$Q_{\mathcal P}(y,\omega)$ for the terminal set under internal seed
$\omega$. Define its optimal worst-case precision-success probability by
\begin{equation}
\begin{aligned}
 S_\tau^{\mathcal C}
 &=\sup_{\mathcal P\in\mathcal C}\inf_{y\in F}
   \Pr_\omega\!\left\{W_{Q_{\mathcal P}(y,\omega)}(y)\leq\tau\right\}\\
 &=1-R_{\mathcal C}^{\phi_\tau},\qquad
   \phi_\tau(w)=\boldsymbol{1}\{w>\tau\},
\end{aligned}
\label{eq:risk_precision_success}
\end{equation}
where $R_{\mathcal C}^{\phi_\tau}$ denotes the corresponding risk in
\eqref{eq:risk_objectives}. The identity follows by complementing the
success event and reversing each supremum or infimum under subtraction
from one; it does not require an optimal policy to attain either bound.
The probability concerns the algorithm's seed at a fixed hidden record,
and equality at the threshold counts as success. Reaching this precision
does not require the interval to exclude zero, so it need not certify a
winning model.

\begin{proposition}[Concavity and the transcript fiber]
\label{prop:risk_concavity_transcript}
For every fixed $Q$, $W_Q$ is nonnegative, finite, concave, and upper
semicontinuous on $F$. For every adaptive execution, its terminal
information fiber is precisely $F_Q(y)$ for its terminal queried set $Q$.
\end{proposition}
\begin{proof}
Choose maximizing and minimizing endpoints $u_k,v_k\in F_Q(y_k)$,
$k=1,2$. For $0\le\lambda\le1$, convexity of $F$ and linearity of the
measurements make $\lambda u_1+(1-\lambda)u_2$ and
$\lambda v_1+(1-\lambda)v_2$ feasible endpoints at
$\lambda y_1+(1-\lambda)y_2$. Their target difference is
$\lambda W_Q(y_1)+(1-\lambda)W_Q(y_2)$, proving concavity.
For upper semicontinuity, take $y_k\to y$ and a subsequence attaining
the limsup of the widths. Compactness supplies a further subsequence
whose endpoint pairs converge to $u,v\in F_Q(y)$. Its limiting
difference is at most $W_Q(y)$.

For the transcript claim, fix the internal random seed. If no query is
made, both fibers are $F$. Otherwise, any
$z\in F_Q(y)$ gives the same answer to the first queried measurement.
Inductively, equal previous answers produce the same next measurement
and the same next answer. The stopping decision is also unchanged.
Thus every completion agreeing with all queried answers follows the
same entire transcript. Conversely, any completion consistent with
that transcript must agree on every queried answer. Query choices and
the stopping time therefore impose no additional restrictions on the
fiber beyond those exact observations. This argument is pointwise in the
seed; no conditioning on a positive-probability seed event is required.
\end{proof}

\subsection{Which width penalties can benefit from adaptation?}
\label{sec:risk_classification}
The next result uses classical minimax theory. The contribution is the
resulting classification for sharp information-fiber widths and its
coordinate-query converse, rather than a new minimax or duality theorem.
The proof has two parts. For concave penalties, a single hardest record
limits every final query set, including those chosen adaptively; Jensen's
inequality supplies a batch policy meeting that limit. For any nonconcave
penalty, a three-coordinate construction realizes a strict midpoint
violation as an adaptive advantage. These arguments establish the universal
claim; the later numerical checks test implementations on particular inputs.

\begin{theorem}[Universal classification of width penalties]
\label{prop:risk_classification}
Let $\phi:[0,\infty)\to\mathbb R$ be finite-valued and nondecreasing.
Randomized batch and randomized adaptive querying have equal minimax
$\phi$-risk for every problem satisfying Section~\ref{sec:risk_setup}
if and only if $\phi$ is concave. More precisely, put
\begin{equation}
 \rho=\min_{\pi\in\Delta(\mathcal Q)}\max_{y\in F}
         \sum_{Q\in\mathcal Q}\pi_Q W_Q(y).
 \label{eq:risk_rho}
\end{equation}
If $\phi$ is concave, then
\begin{equation}
 R_{\mathrm{adaptive}}^{\phi}
 =R_{\mathrm{batch}}^{\phi}=\phi(\rho).
 \label{eq:risk_concave_value}
\end{equation}
Every batch policy attaining~\eqref{eq:risk_rho} simultaneously attains
this equality for every such concave $\phi$. If $\phi$ is not concave,
a strict adaptive improvement already occurs on a three-coordinate
compact convex polytope, with all coordinates queryable at unit cost
and a hard budget of two.
\end{theorem}
\begin{proof}
The function $\sum_Q\pi_QW_Q(y)$ is affine and continuous in $\pi$,
and concave and upper semicontinuous in $y$. Sion's minimax theorem \citep{sion1958}
therefore yields
\begin{equation}
 \rho=\max_{y\in F}\min_{Q\in\mathcal Q}W_Q(y).
 \label{eq:risk_hardest_record}
\end{equation}
All values are finite. Writing $G(\pi)=\max_{y\in F}\sum_Q\pi_QW_Q(y)$,
the bound $|G(\pi)-G(\pi')|\le D_T\|\pi-\pi'\|_1$ proves continuity.
Thus the minimum in~\eqref{eq:risk_rho} is attained on the compact simplex.
The maximum in~\eqref{eq:risk_hardest_record} is attained because the
minimum of finitely many upper semicontinuous functions is upper
semicontinuous on compact $F$. Hence there is one record $y^\star\in F$
such that $W_Q(y^\star)\ge\rho$ for every admissible $Q$.
By Proposition~\ref{prop:risk_concavity_transcript}, every execution of
every adaptive policy has terminal width at least $\rho$ at this same
record. Monotonicity gives the lower bound $\phi(\rho)$.

Let $\pi^\star$ attain~\eqref{eq:risk_rho}. Concavity, Jensen's
inequality, and monotonicity give, for every $y\in F$,
\[
 \sum_Q\pi^\star_Q\phi(W_Q(y))
 \le \phi\!\left(\sum_Q\pi^\star_QW_Q(y)\right)
 \le\phi(\rho).
\]
This proves the upper bound and simultaneous optimality. No continuity
of $\phi$ at zero is needed; the finite-mixture form of Jensen's
inequality suffices. For example, $\phi(w)=\boldsymbol{1}\{w>0\}$ is
concave on $[0,\infty)$ despite its jump at zero.

For necessity, a nonconcave, nondecreasing finite-valued $\phi$ must
have a strict midpoint-concavity violation: there exist $0\le l<h$ such
that, with $m=(l+h)/2$,
\begin{equation}
 \phi(m)<\{\phi(l)+\phi(h)\}/2.
 \label{eq:risk_midpoint_violation}
\end{equation}
Indeed, if midpoint concavity held, for $0\le x\le z$ and dyadic
$q_k\uparrow\theta$ its dyadic inequalities and monotonicity would give
\[
\phi((1-\theta)x+\theta z)
\ge\phi((1-q_k)x+q_kz)
\ge(1-q_k)\phi(x)+q_k\phi(z).
\]
Taking the limit only on the right proves full concavity without a
continuity assumption, a contradiction. Set $d_0=h-l$ and consider
\begin{equation}
 F_{l,h}=\{(a,b,c):0\le a\le1,\ 0\le b\le l+d_0a,
                  \ 0\le c\le h-d_0a\},\qquad T=ha+b+c.
 \label{eq:risk_converse_polytope}
\end{equation}
All three coordinates are available at unit cost, with budget two.
The widths after querying $\{a,b\}$ and $\{a,c\}$ are respectively
$h-d_0a$ and $l+d_0a$. An adaptive policy first observes $a$ and then
queries the coordinate with the larger remaining capacity. Its terminal
width is at most $m$, so its risk is at most $\phi(m)$.
At $(a,b,c)=(1/2,0,0)$, these two widths equal $m$, while the width
after querying $\{b,c\}$ is $h$. Thus every admissible terminal set
has width at least $m$, and the adaptive optimum is exactly $\phi(m)$.

Now put equal adversarial mass on $(0,0,0)$ and $(1,0,0)$.
Each batch $\{a,b\}$ or $\{a,c\}$ has average penalty
$\{\phi(l)+\phi(h)\}/2$. The batch $\{b,c\}$ has width $h$ at both
records, so its average penalty is no smaller. A batch using fewer
coordinates cannot improve upon an extension to two coordinates,
because more observations only shrink the fiber. Every randomized
batch policy consequently has worst-case risk at least
$\{\phi(l)+\phi(h)\}/2>\phi(m)$. This is the required strict separation.
\end{proof}

The universal quantifier is essential. A nonconcave penalty permits an
adaptive advantage somewhere in this class; it need not yield an
advantage for every feasible set. The theorem covers heterogeneous fixed
costs and overlapping constraints at the level of information theory.
It does not imply a polynomial-time optimizer for an arbitrary feasible
set or budget.

\subsubsection{A three-coordinate example}
\label{sec:risk_small_example}
Take $0\le a\le1$, $0\le b\le a$, $0\le c\le1-a$, and
$T=a+b+c$, with two unit-cost queries. The widths for batches
$\{a,b\}$, $\{a,c\}$, and $\{b,c\}$ are $1-a$, $a$, and $1-b-c$.
Querying $a$ first and then $c$ if $a\le1/2$, or $b$ otherwise, leaves
width $\min(a,1-a)\le1/2$.

\begin{table}[htbp]
\centering
\caption{Exact constructed-example values. These are mathematical
guarantees, not forecasting improvements measured on public data.}
\label{tab:risk_three_coordinate}
\begin{tabular}{lcc}
\toprule
Criterion & Best randomized batch & Best adaptive policy\\
\midrule
Worst expected width & $1/2$ & $1/2$\\
Worst expected squared width & $1/2$ & $1/4$\\
Worst-case success, width $\le3/4$ & $1/2$ & $1$\\
\bottomrule
\end{tabular}
\end{table}

An equal mixture of $\{a,b\}$ and $\{a,c\}$ attains the batch values.
The endpoint records $(0,0,0)$ and $(1,0,0)$ give the corresponding
batch lower bounds. At $(1/2,0,0)$ every admissible terminal set has
width at least $1/2$, establishing the adaptive moment lower bounds.
The precision separation holds for every threshold $1/2\le\tau<1$.
Reaching a width threshold is distinct from certifying the sign of $T$:
a narrow sharp interval may still contain zero.

\subsection{Growing moment and precision separations}
\label{sec:risk_growing}
For $K\ge2$, let $a_1,\ldots,a_K\ge0$ sum to one and impose
$0\le b_j\le a_j$. The actual queryable coordinates are
$a_1,\ldots,a_{K-1},b_1,\ldots,b_K$; $a_K$ is determined by the
first $K-1$ capacities. Define
\begin{equation}
 T=\sum_{i<K}a_i+\sum_{j=1}^K b_j,\qquad B=2K-2.
 \label{eq:risk_K_setup}
\end{equation}
This is a compact convex polytope in $2K-1$ dimensions, and every
coordinate has the same unit cost.

\begin{theorem}[Sharp growing separations]
\label{prop:risk_K_family}
For~\eqref{eq:risk_K_setup}, the optimal worst expected width is $1/K$
for both randomized batch and adaptive policies. For every $p>1$,
\begin{equation}
 R_{\mathrm{batch}}^{w^p}=1/K,\qquad
 R_{\mathrm{adaptive}}^{w^p}=1/K^p.
 \label{eq:risk_K_moments}
\end{equation}
For required width $\tau=1/K$, the best worst-case success probabilities
are $1/K$ for randomized batches and $1$ for adaptive policies.
\end{theorem}
\begin{proof}
It suffices to consider batches using the full budget: adding observations
cannot increase width. Such a batch omits exactly one coordinate.
Omitting $b_j$ leaves width $a_j$; omitting $a_i$, $i<K$, leaves width
$a_i+a_K-b_i-b_K$. The latter follows because the remaining feasible
range of $a_i$ is $[b_i,a_i+a_K-b_K]$ and its coefficient in $T$ is one.
An adaptive policy queries all explicit capacities and then all $b_j$
except one having the smallest capacity. Its width is
$\min_j a_j\le1/K$. At the record with all $a_j=1/K$ and all $b_j=0$,
every full-budget batch has width at least $1/K$. This proves the
adaptive lower bounds and the mean-width lower bound.

The randomized batch policy omitting each $b_j$ with probability $1/K$
has expected width exactly $1/K$. Its expected $p$th power is
$K^{-1}\sum_j a_j^p\le1/K$, since $0\le a_j\le1$ and $p>1$.
For the reverse moment bound, average over the $K$ records
$a=e_j$, $b=0$. A batch omitting a $b_j$ has average $p$th-power width
$1/K$; one omitting an explicit $a_i$ has average $2/K$.
Therefore every batch distribution has worst-case moment at least $1/K$.

For precision, omitting a uniformly random $b_j$ succeeds with
probability at least $1/K$ at every record, because at least one capacity
is no greater than $1/K$. To prove the matching upper bound, use $K$
records indexed by $j$, with $b=0$,
\[
 a_j=\frac{1}{2K},\qquad
 a_i=\frac{2K-1}{2K(K-1)}\quad(i\ne j).
\]
These capacities sum to one. Only the batch omitting $b_j$ reaches
width at most $1/K$ at record $j$: every other $b$-omitting batch has larger
width, and every $a$-omitting batch has width $a_i+a_K>1/K$.
If $q_j$ is the probability of omitting $b_j$, worst-case success is
at most $\min_j q_j\le1/K$. The adaptive policy always succeeds.
The precision upper bound uses these interior capacity records;
checking only vertices would not establish it.
\end{proof}

The moment ratio is $K^{p-1}$, while the ratio of worst-case success
probabilities at precision $1/K$ is $K$. These separations use a
structured constructed feasible set and a budget querying all but one
coordinate. They establish the possible difference between objectives;
they do not establish comparable gains at small relative budgets in a
deployment or in the historical forecasting study.
\subsection{Computing policies and hardest records on polytopes}\label{sec:general-algorithms}
Let $F=\{y\in\mathbb R^n:Ay\le b\}$ be nonempty and compact, and let
$T(y)=c^\top y+d$. Equalities can be represented by paired inequalities.
The constant $d$ does not affect width. The first formulation allows any
declared finite family $\mathcal Q$ of legal batches; $L_Q$ denotes the matrix
of their exact linear measurements.

\subsubsection{An enumerative hardest-record formulation}
Introduce a common record $y$, endpoint witnesses $u_Q,v_Q$ for each
$Q\in\mathcal Q$, and a scalar $\eta$. Consider
\begin{equation}\label{eq:general-hardest-lp}
\begin{aligned}
\max\quad &\eta\\
\text{subject to}\quad &Ay\le b,\quad Au_Q\le b,\quad Av_Q\le b
                         &&(Q\in\mathcal Q),\\
&L_Qu_Q=L_Qy,\quad L_Qv_Q=L_Qy &&(Q\in\mathcal Q),\\
&\eta\le c^\top(u_Q-v_Q) &&(Q\in\mathcal Q).
\end{aligned}
\end{equation}
For fixed $y$, the independent witnesses for batch $Q$ can attain its sharp
width $W_Q(y)$. Thus the largest feasible $\eta$ at that record is
$\min_Q W_Q(y)$, and the optimum of \eqref{eq:general-hardest-lp} is $\rho$.
Its common record is a certificate that every legal batch leaves width at
least $\rho$ there.

The nonnegative dual multipliers $\pi_Q$ of the last constraints form an
optimal randomized batch policy. Indeed, the free variable $\eta$ forces
$\sum_Q\pi_Q=1$. Dualizing these constraints first yields
\[
\min_{\pi\in\Delta(\mathcal Q)}
 \max_{y,u_Q,v_Q}\sum_Q\pi_Qc^\top(u_Q-v_Q)
 =\min_{\pi\in\Delta(\mathcal Q)}\max_{y\in F}
       \sum_Q\pi_Q W_Q(y)=\rho,
\]
where the endpoint feasibility constraints remain in the inner maximization.
Ordinary linear-programming strong duality supplies the equality and the
policy; no new duality theorem is asserted. In the implementation, a separate
LP fixes the extracted $\pi$ and maximizes its expected width, checking the
upper bound independently of the constraints involving $\eta$.

With $M=|\mathcal Q|$ batches, the formulation has $n(1+2M)+1$ variables.
It is polynomial in the explicitly enumerated batch family, which may be
exponential in the number of available measurements. It is a small-instance
benchmark and certificate generator; the aggregate-specific compact
formulation avoids this enumeration in its structured setting.

\subsubsection{An exact deterministic two-query algorithm}\label{sec:two-query}
Assume every one of the $n\ge2$ coordinates is queryable at unit cost and the
hard budget is two. An unused query can be added without increasing width,
so considering two distinct coordinates includes policies that stop earlier.
For a first query $i$ and feasible answer $t$, define
\begin{equation}\label{eq:conditional-two-query}
g_{ij}(t)=\max\{c^\top(u-v):Au\le b,\ Av\le b,\
                u_i=v_i=t,\ u_j=v_j\},\qquad j\ne i.
\end{equation}
This maximizes the terminal width over all feasible records having first
answer $t$, after second query $j$. The best deterministic policy beginning
with $i$ therefore has worst width
$\rho_i=\max_t\min_{j\ne i}g_{ij}(t)$.

\begin{proposition}\label{prop:two-query-lp}
For each $i$, $\rho_i$ is the value of the following single linear program,
with one endpoint pair $(u^j,v^j)$ for each $j\ne i$:
\begin{equation}\label{eq:two-query-lp}
\begin{aligned}
\max\quad &\eta\\
\text{subject to}\quad &Au^j\le b,\quad Av^j\le b &&(j\ne i),\\
&(u^j)_i=(v^j)_i=t,\quad (u^j)_j=(v^j)_j &&(j\ne i),\\
&\eta\le c^\top(u^j-v^j) &&(j\ne i).
\end{aligned}
\end{equation}
The optimal deterministic adaptive two-query worst-width value is
$\min_i\rho_i$.
\end{proposition}
\begin{proof}
At fixed feasible $t$, each endpoint pair has only its own constraints and
can attain $g_{ij}(t)$. All inequalities on $\eta$ are feasible precisely up
to their smallest attainable upper bound, $\min_{j\ne i}g_{ij}(t)$.
Maximizing over the shared $t$ proves the first claim.
For any deterministic policy beginning with $i$, its chosen second coordinate
at answer $t$ has conditional worst width at least $\min_jg_{ij}(t)$.
Conversely, choosing a minimizing $j$ after observing $t$ attains that value.
Finite tie breaking gives a measurable selection because these polyhedral
value functions are piecewise linear on their feasible interval. Finally,
minimize over the first coordinate.
\end{proof}

For rational input, $n$ LPs of polynomial description size choose the first
coordinate; after its answer, $n-1$ conditional LPs implement the second
choice. This is a polynomial-time deterministic two-query procedure in the
standard bit-complexity model, including the encoding length of rational
query answers. For arbitrary exact answers, the policy is still defined by
the conditional LPs; the polynomial bit-complexity statement applies to
rationally encoded inputs and answers. It does not solve arbitrary-budget
randomized precision-risk optimization or minimize expected requests until
certification. No grid of possible first answers is used.
\subsubsection{A precision consequence of sparse-policy bounds}
% Research draft: AI-assisted derivation and internal mathematical review.
% Not a human-specialist validation or historical-priority claim.
The dimension bound below adapts established sparse-policy reasoning.
\citet[Theorem~2 and Example~5]{kurtz2026} gives the corresponding
dimension-plus-one policy bound for continuous concave scenario costs,
including conditional information discovery. Conditional terminal widths
fit that framework on polytopes. The direct Helly proof below also allows
compact convex $F$, for which these widths are upper semicontinuous but
need not be continuous, and translates the sparse cover into precision
success under uniform randomization.
\begin{theorem}[A dimension-dependent batch guarantee, adapted from sparse-policy bounds]
\label{thm:helly-precision}
Let $F$ be nonempty, compact and convex, and let $T(x)=c^\top x+d$.
Fix a first-stage exact linear observation $L_0x$, whose feasible-answer
set $I=L_0F$ has affine dimension $r$. After this answer, a deterministic
policy chooses one of finitely many second-stage linear measurement
batches $L_1,\ldots,L_m$. Every combined batch $(L_0,L_j)$ must obey the
same declared hard budget, and the eligible second-stage batches must not
depend on the first answer. Suppose the policy's terminal sharp interval
has width at most $\tau$ for every $x\in F$.
Then some randomized nonadaptive policy supported on at most $r+1$
of these combined batches has worst-case precision success at least
$1/(r+1)$.
In particular, a sure-precision policy using two scalar queries yields
a randomized-batch success guarantee of at least $1/2$, even when its
conditional worst widths are not affine functions of the first answer.
\end{theorem}
\begin{proof}
For $t\in I$, define the conditional worst terminal width
\[
 g_j(t)=\max\{c^\top(u-v):u,v\in F,\ L_0u=L_0v=t,
                                      \ L_ju=L_jv\}.
\]
The feasible set is nonempty and compact, so the maximum is attained.
The expression is exactly the largest terminal sharp width among records
with first answer $t$: an endpoint pair with equal first and second
answers belongs to a common terminal fiber, and the endpoints of every
such fiber are feasible in this maximization.

Each $g_j$ is concave on $I$. Indeed, convex combinations of optimizing
endpoint pairs at $t_1,t_2$ are feasible at the corresponding convex
combination of answers, and attain the same convex combination of their
objective values. The function is also upper semicontinuous: for a
convergent sequence of answers, compactness gives convergent subsequences
of optimizing endpoints, whose limits are feasible at the limiting
answer. Consequently, $g_j$ is finite and measurable.

For a fixed first answer, the policy chooses a single second-stage batch;
its sure guarantee therefore implies $\min_jg_j(t)\leq\tau$ throughout
$I$. The strict failure sets
\[
 H_j=\{t\in I:g_j(t)>\tau\}
\]
are convex by concavity, and $\bigcap_{j=1}^m H_j=\varnothing$.
Finite Helly's theorem \citep[Theorem~1.1]{amenta2016}, applied in the affine hull of $I$, gives a
subfamily $J$ of at most $r+1$ sets with empty intersection. This finite
form of Helly's theorem does not require the sets to be closed; thus
using strict failure sets preserves success at equality with $\tau$.

Choose uniformly among the combined batches indexed by $J$. For every
fixed record $x$, its answer $t=L_0x$ lies outside at least one $H_j$,
so that batch's actual terminal width is at most $g_j(t)\leq\tau$.
Its precision success probability is therefore at least $1/|J|$, and
hence at least $1/(r+1)$, for every record.
\end{proof}


The empty second-stage action includes stopping after the first stage.
Fixed unequal costs are allowed if all combined batches obey the budget;
the action dictionary and latent record are fixed before randomization.
The $K$-capacity construction in Theorem~\ref{prop:risk_K_family} is sharp
with first-answer dimension $r=K-1$ and optimal batch success $1/K$.
For scalar first answers, at most two second-stage actions suffice to
retain the optimal sure-width guarantee. On polytopes, test each singleton
or pair of actions with the shared-first-answer LP in
Proposition~\ref{prop:two-query-lp}; there are at most
$\binom m2+m$ candidates. This sparse policy need not minimize conditional
width at every individual answer. Helly-based sparsification for
convex-function optimization is classical
\citep[Section~2, Theorem~1]{rockafellar1965}. The contribution here is
its certificate-risk interpretation and sharpness through admissible
information queries, rather than a new general sparse-policy theorem.
\subsubsection{Adaptive precision from cumulative-counter queries}
\label{sec:counter-precision-method}
Consider a six-hour block with nonnegative hourly increments
$z=(z_1,\ldots,z_6)$ and retained cumulative registers $C_0,C_6$,
all expressed in the same energy units.
Writing $E=C_6-C_0$, the initial feasible set is
$F_E=\{z\geq0:\sum_{h=1}^6z_h=E\}$.
For two frozen forecast vectors $a,b$, their block MSE contrast is
\[
 D(z)=d+c^\top z,\qquad
 d=\frac16\sum_{h=1}^6(a_h^2-b_h^2),\quad
 c_h=\frac{2(b_h-a_h)}6.
\]
The permitted query at boundary $k\in\{1,\ldots,5\}$ returns
$C_k$, equivalently the prefix sum $s_k=C_k-C_0$.
The budget is at most two distinct internal registers per pair and block.
This is a declared cumulative-register access interface, rather than the
coordinate-increment interface of the earlier recovery experiment.
The register queries are exact linear measurements covered by the general
information model. Adaptive scalar certification also appears in monotone
integral estimation \citep[Section~4.1]{gaillard2023}; here the target is an
affine contrast of interval increments under a fixed hard budget.

For integer boundaries $0\leq p<q\leq6$, define
$r(p,q)=\max_{p<h\leq q}c_h-\min_{p<h\leq q}c_h$.
With queried boundaries $Q$, let $0=q_0<\cdots<q_m=6$ include
the outer boundaries and $Q$, and put $s_0=0$, $s_6=E$.
Nonnegativity and the segment totals give the sharp width
\begin{equation}
 W_Q(z)=\sum_{\ell=1}^{m}
 (s_{q_\ell}-s_{q_{\ell-1}})r(q_{\ell-1},q_\ell).
 \label{eq:counter-width}
\end{equation}
Both endpoints are attainable by placing each segment's mass at a
minimum- or maximum-coefficient hour. The initial width is $W_0=Er(0,6)$.
Tasks with $W_0=0$ are already precise and require no queries.

\begin{proposition}[Two best-side splits for an arbitrary block]\label{prop:counter-two-query}
Let $n\geq3$, $E>0$, $x\geq0$, $\sum_{h=1}^n x_h=E$, and let the target be
$d+\sum_hc_hx_h$. Each query returns one internal prefix sum. Write
$r(p,q)=\max_{p<h\leq q}c_h-\min_{p<h\leq q}c_h$.
For first boundary $k$, put $\ell_k=r(0,k)$ and $r_k=r(k,n)$.
Define each of the following quantities whenever its corresponding
side has an eligible second boundary:
\[
 \alpha_k=\min_{1\leq j<k}\max\{r(0,j),r(j,k)\},\qquad
 \beta_k=\min_{k<j<n}\max\{r(k,j),r(j,n)\}.
\]
Choose a minimizing boundary on each side, breaking ties by smallest index.
When both sides are nonempty, for normalized first answer $t=s_k/E$,
the optimal conditional worst width is
\[
 E\min\{\alpha_k t+r_k(1-t),\ \ell_k t+\beta_k(1-t)\}.
\]
Thus only one best boundary on each side is needed, independently of $t$.
If $\Delta_k=\ell_k+r_k-\alpha_k-\beta_k>0$, the switch is
\[
 t_k^\star=\frac{r_k-\beta_k}{\Delta_k}\in[0,1],
\]
and the worst width for this first boundary is
\[
 A_k=E\max\left\{\alpha_k,\beta_k,
          \frac{\ell_kr_k-\alpha_k\beta_k}{\Delta_k}\right\}.
\]
Use the best right split below the switch and the best left split above
it; either is valid at equality. If $\Delta_k=0$, the two lines coincide
and $A_k=E\max\{\ell_k,r_k\}$. At $k=1$ only the right split exists and
$A_1=E\beta_1$; at $k=n-1$ only the left split exists and
$A_{n-1}=E\alpha_{n-1}$. The optimal deterministic two-query guarantee is
\begin{equation} A=\min_{1\leq k<n} A_k.\label{eq:counter-adaptive-value}\end{equation}
\end{proposition}
\begin{proof}
All second queries left of $k$ have the common right-side contribution
$r_k(1-t)$; minimizing their remaining term $t$ times a split range
selects the same minimizing boundary for every $t\geq0$. The right-side
argument is identical. Moreover $\alpha_k\leq\ell_k$ and
$\beta_k\leq r_k$. Hence $\Delta_k\geq0$, and when it is positive the
displayed intersection belongs to $[0,1]$. The maximum of the lower
envelope occurs at an endpoint or that intersection. Its endpoint
values are $\beta_k,\alpha_k$, and its intersection value is the displayed
quotient. If $\Delta_k=0$, both nonnegative differences vanish, so the
lines coincide. At an outer eligible boundary the single-hour side has
zero coefficient range, giving the stated one-line values.
\end{proof}

All contiguous ranges can be built in $O(n^2)$ arithmetic operations
and stored in $O(n^2)$ space. Scanning the candidates for every
$\alpha_k,\beta_k$ costs another $O(n^2)$ operations; the closed forms
then cost $O(n)$. This is an arithmetic-operation bound for the exact-total,
uncapped two-prefix problem, not a bit-complexity or noisy-polytope claim.
For $n=2$, the single internal register determines both increments;
for $n=1$ the total already determines the only increment.


\paragraph{An upper certificate for randomized-batch precision.}
Returning to the six-hour case, there are ten two-boundary batches.
Let $z^{(h)}=Ee_h$ be the six feasible
single-hour mass profiles, and define
$M_{hQ}(\tau)=\mathbf1\{W_Q(z^{(h)})\leq\tau\}$.
The finite-profile program
\begin{equation}
 U(\tau)=\min_{w\geq0,\,\sum_hw_h=1,\,u}u
 \quad\text{subject to}\quad
 \sum_hw_hM_{hQ}(\tau)\leq u\quad\text{for every batch }Q
 \label{eq:counter-precision-upper}
\end{equation}
upper-bounds the worst-case precision success of \emph{every} randomized
batch policy. Indeed, its minimum success over $F_E$ is at most its success
averaged over these profiles under $w$, which is at most
$\max_Q\sum_hw_hM_{hQ}$. This remains valid for batches with fewer queries:
adding an unused register cannot increase width. The exported $w$ therefore
provides a checkable upper certificate. The finite profile set need not
contain a worst completion for an indicator objective, so $U(\tau)$ is not
asserted to equal the full precision optimum in general. Below, write
$S_B^\star=S_\tau^{\mathrm{batch}}$ and
$S_A=S_\tau^{\mathrm{adaptive}}$ for the precision-success values defined
in \eqref{eq:risk_precision_success}, at this fixed instance and threshold.
If $A\leq\tau$ and $U(\tau)<1$, the adaptive policy succeeds for every
completion, whereas no randomized batch can do so. In floating-point
classification, batch widths within $10^{-9}W_0$ above the threshold count
as success, and adaptive certification requires a margin of
$10^{-9}W_0$; exact constructed boundary cases are treated analytically.

\paragraph{An exact counter example.}
Take $E=1$ and $c=(0,0,1,1,0,1)$, so $W_0=1$.
Query boundary $4$ first, obtaining $t=s_4$.
Batch $\{2,4\}$ has width $1-t$, while batch $\{4,5\}$ has width $t$.
Querying $5$ when $t<1/2$ and $2$ otherwise thus guarantees
width $\min\{t,1-t\}\leq1/2$.
Conversely, at $z=e_5$ only batch $\{4,5\}$ attains width below $1$;
that batch has width $1$ at $z=e_2$. No batch succeeds at both profiles
for any threshold $\tau<1$. Weights $w_2=w_5=1/2$ in
\eqref{eq:counter-precision-upper} give $U\leq1/2$.
An equal mixture of $\{2,4\}$ and $\{4,5\}$ succeeds with probability
at least $1/2$ on every completion when $\tau\geq1/2$.
Consequently, for every $\tau\in[1/2,1)$ the optimal adaptive worst-case
success is $1$ and the optimal randomized-batch worst-case success is
$1/2$. These are exact mathematical values, not measured application
performance or acquisition-cost savings.

\begin{corollary}[The two-register half-success floor]
\label{prop:counter-half-floor}
If a deterministic two-register policy guarantees width at most $\tau$,
a randomized batch supported on at most two batches has worst-case
precision success at least $1/2$. Thus an adverse-profile upper bound
$U(\tau)=1/2$ proves the exact batch optimum for such an instance.
\end{corollary}
\begin{proof}
Apply Theorem~\ref{thm:helly-precision} with scalar first answer.
For the exact-total prefix model the conditional functions are affine:
their good sets are prefixes or suffixes of $[0,1]$. The longest good
prefix and earliest good suffix cover all answers; mix their batches
equally, breaking geometric ties by smallest second-query index.
\end{proof}
% Independently reviewed operational corollary; no new historical-priority claim.
\begin{corollary}[Deciding whether feedback is necessary]
\label{cor:counter-feedback-triage}
In the cumulative-counter model, let $\mathcal B$ contain the ten
two-register batches, and define
\[
 D=\min_{Q\in\mathcal B}\max_{z\in F_E}W_Q(z)
   =\min_{Q\in\mathcal B}\max_{1\leq h\leq6}W_Q(Ee_h).
\]
Let $A$ be the exact deterministic adaptive value in
\eqref{eq:counter-adaptive-value}, and let $U(\tau)$ be the
finite-profile upper certificate in \eqref{eq:counter-precision-upper}.
For every $\tau\geq0$,
\[
 D\leq\tau
 \quad\Longleftrightarrow\quad U(\tau)=1
 \quad\Longleftrightarrow\quad S_B^\star=1.
\]
Consequently, the available forecasts and outer registers determine
three mutually exclusive decisions before any optional retrieval:
\begin{enumerate}
\item If $D\leq\tau$, a fixed batch guarantees the requested precision.
\item If $A\leq\tau<D$, feedback enables a deterministic guarantee that
      every randomized batch fails to match.
\item If $\tau<A$, no deterministic adaptive policy with two register
      queries guarantees the requested precision.
\end{enumerate}
In the second case,
$1/2\leq S_B^\star\leq U(\tau)\leq5/6$.
The lower bound uses Corollary~\ref{prop:counter-half-floor}; when
$U(\tau)=1/2$, the randomized-batch optimum is exactly $1/2$.
\end{corollary}
\begin{proof}
For a fixed batch, \eqref{eq:counter-width} is a linear function of
$z$. It therefore attains its maximum over the energy simplex at a
concentration profile $Ee_h$, which proves the expression for $D$.
If $D\leq\tau$, one batch succeeds at all six profiles and, by
linearity, at every completion. Thus $S_B^\star=1$ and $U(\tau)=1$.
If $D>\tau$, every batch fails at at least one profile. The uniform
distribution on the six profiles consequently assigns every batch
success probability at most $5/6$. It is a feasible adverse distribution
in \eqref{eq:counter-precision-upper}, so
$S_B^\star\leq U(\tau)\leq5/6<1$. This proves both equivalences.
Adaptation includes fixed batches, hence $A\leq D$. The three cases
then follow from the defining optimality of $A$ and $D$.
\end{proof}

This decision rule uses only the forecast coefficients, retained total,
query dictionary and requested precision. It identifies when a sequential
query response is necessary for a guarantee, without measuring retrieval
latency, financial cost or the frequency of adverse energy profiles.
The finite-profile program is decisive for whether batch success equals
one; it remains only an upper bound on intermediate success values unless
a matching guarantee is supplied.
\section{Results}\label{sec:evaluation}
The application asks which admissible retrieval method meets a required
precision. It separates information available before a query, guarantees
over all compatible records, replay on recorded answers, and actual API
measurements. All policies in each comparison share the same initial
information and measurement model. Scalar readings, network requests and
feedback rounds are different resources. The protocol must identify which
resource is limited before comparing methods.
\subsection{When two counter readings need feedback}
\label{sec:counter-precision-results}
The first replay uses three actual cumulative consumption counters \citep{blem}, 91
eligible days per meter, and four frozen forecasters (persistence, ridge,
boosted trees and an MLP). Each of their six pairs is audited on six-hour
blocks ending at 06:00, 12:00 and 18:00 UTC: 4,914 tasks in total.
The auditor receives the block's endpoint counters and day-start forecasts.
Next-day persistence copies are unavailable at these intraday cutoffs under
the declared issue schedule. No production arrival times are inferred.
Source and fitting details are in Online Resource~1, S3.

The restricted interface offers five internal hourly counters and a budget
of two scalar readings per standalone pair/block audit. This is a specified
access replay, not a provider tariff or a demonstrated deployed service.
Hourly-increment queries and bulk exports are separate interfaces, examined
below. Every policy uses the same retained total, nonnegativity constraints
and forecast-derived coefficients. The primary target, fixed before the
geometry and answer replay, is $W\leq W_0/2$. This is a precision requirement;
it need not establish which model wins. Query choices use only permitted
information and previous queried answers. Unqueried answers enter outcome
evaluation after the choice. The shared preparation files are not physically
isolated, so this is an audited dependency rule rather than a security boundary.

Table~\ref{tab:counter-guarantees} gives the complete decision partition.
A fixed batch suffices on 88.95\% of tasks. Feedback changes the attainable
deterministic guarantee on 5.90\%, spanning all three meters and all six
model pairs. The remaining 5.15\% require more information, a larger budget
or a weaker requirement. The first 20 days per meter contributed 1,080 tasks;
the unchanged continuation contributed 3,834. These reused meters are not
an untouched-source replication.

\begin{table}[htbp]\centering\small
\caption{Decisions before optional answers, at $W\leq W_0/2$ with two
readings. ``Feedback needed'' means $A\leq\tau<D$; ``unresolved'' means
$\tau<A$ for deterministic policies. All 290 feedback cases have exact
optimal randomized-batch worst-case success $1/2$.}
\label{tab:counter-guarantees}
\begin{tabular}{lrrrr}\toprule
Meter & Tasks & Fixed batch suffices & Feedback needed & Unresolved\\\midrule
1 & 1,638 & 1,460 & 97 & 81\\
2 & 1,638 & 1,430 & 103 & 105\\
3 & 1,638 & 1,481 & 90 & 67\\
All & 4,914 & 4,371 & 290 & 253\\\bottomrule
\end{tabular}\end{table}

For every feedback case, a constructed equal mixture of two batches
succeeds on at least half of all random choices at every feasible record.
A finite adverse-profile certificate gives the reverse bound for every
randomized batch. These establish a continuum optimum, not just a sampled
upper bound: $S_B^\star\geq1/2$ and $S_B^\star\leq1/2$.
Consequently, the minimum, quartiles, median and maximum of the 290 exact
bounds all equal $1/2$; none is below 40\% or 30\%.
Corollary~\ref{prop:counter-half-floor} explains this structural floor.
The worked example below supplies the complete argument for one recorded
comparison. S6 retains the earliest qualifying witness from each meter,
including exact coefficients, covers and adverse weights.

\paragraph{One fully inspectable recorded comparison.}\label{sec:worked-counter-example}
Take meter~1, persistence versus boosting, 1 October 2017, 06:00--12:00
UTC, the earliest qualifying case for that meter. Its retained total is
$E=4.6129961798$~kWh. With $z=y/E$ and
$q=(c-\min c)/(\max c-\min c)$, we have $z\geq0$, $\sum_hz_h=1$ and
$W/W_0$ equal to the width of $q^\top z$. Exact conversion of the saved
coefficients gives the following outward enclosures; every displayed
endpoint is an exact decimal rational:
\[
q\in\{1\}\times[.9061,.9062]\times\{0\}
\times[.5717,.5718]\times[.8607,.8608]\times[.2988,.2989].
\]
The initial squared-loss width satisfies
$1.31295<W_0=E(\max c-\min c)<1.31296$; the required width is
$\tau=W_0/2$ (approximately $0.656478$).

Query boundary~3 first and observe $t=z_1+z_2+z_3$.
Query boundary~5 if $t\leq1/4$, and boundary~2 otherwise.
The coefficient enclosures alone give the conditional bounds
\[
\begin{array}{ll}
t+.2891(1-t)\leq .466825,&0\leq t\leq1/4,\\
.0939t+.5620(1-t)\leq .444975,&1/4\leq t\leq1.
\end{array}
\]
Both are strictly below $1/2$, proving success for every feasible record.

For the matching batch limitation, place adverse probability $1/2$ each on
$z=e_2$ and $z=e_5$. The table lists success at the same half-width threshold
for all ten batches ($ij=\{i,j\}$); the enclosures certify every entry.
\begin{center}\small
\begin{tabular}{lrrrrrrrrrr}\toprule
Batch &12&13&14&15&23&24&25&34&35&45\\\midrule
$e_2$ &1&0&0&0&1&1&1&0&0&0\\
$e_5$ &0&0&0&0&0&0&0&0&1&1\\\bottomrule
\end{tabular}
\end{center}
Every column averages at most $1/2$, limiting every randomized batch.
Conversely, an equal mixture of $\{3,5\}$ and $\{2,3\}$ succeeds on
at least one branch for every $t$, giving the continuum lower bound $1/2$.
Thus $S_A=1$ and $S_B^\star=1/2$. Smaller batches cannot improve precision.
Online Resource~1 retains the exact coefficients and an independent
rational checker for this printed witness. This is the existing exact-total
case; it is not a new sample or a measurement-error guarantee.


The replay compares exact adaptive worst-width planning, the best fixed
worst-width batch, a mean-width-optimal randomized batch, uniform batching
and greedy adaptation. Greedy's first-query ties use the best two-query
guarantee; its second query minimizes conditional worst width. Batch ties
minimize uniform-vertex mean width. The mean-optimal policy uses that same
answer-independent secondary criterion on its optimal face. Randomized
outcomes are exact expectations over ten batch probabilities. The added
precision-optimal cover is defined only on the 290-case subset and uses
coefficient-based, answer-independent tie breaking.

\begin{table}[htbp]\centering\small\setlength{\tabcolsep}{4pt}
\caption{Recorded replay, two readings per audit. P: precision success;
W: mean width relative to $W_0$; O: interval excludes zero. All entries
are percentages. The subset is selected by pre-answer guarantees.}
\label{tab:counter-replay}
\begin{tabular}{lrrrrrr}\toprule
 & \multicolumn{3}{c}{All 4,914 tasks} & \multicolumn{3}{c}{290-case subset}\\
Policy & P & W & O & P & W & O\\\midrule
Exact adaptive & 98.64 & 18.59 & 50.49 & 100.00 & 33.72 & 28.97\\
Best deterministic batch & 98.23 & 18.87 & 49.84 & 97.24 & 35.28 & 25.52\\
Mean-optimal random batch & 92.68 & 22.83 & 45.42 & 80.01 & 41.21 & 24.43\\
Precision-optimal cover & -- & -- & -- & 75.34 & 43.12 & 22.76\\
Uniform random batch & 63.04 & 42.59 & 24.79 & 41.79 & 53.10 & 18.76\\
Greedy adaptive & 97.96 & 18.84 & 50.33 & 96.21 & 34.67 & 28.62\\\bottomrule
\end{tabular}\end{table}
\FloatBarrier

The large worst-case success gap is not a large average replay gain.
Relative to the best deterministic batch, adaptation gains 0.407 percentage
points over all tasks: 30 are adaptive-only successes and ten are batch-only
successes. On the subset, eight are adaptive-only and none batch-only.
The precision-optimal cover can perform worse on recorded answers than
policies optimizing other risks. The prespecified 10\% prevalence target
was not met (5.90\%); it is not replaced by a secondary threshold.
All five predeclared precision fractions, paired ordering outcomes and
per-meter summaries remain in S6. Blocks and model pairs are dependent.

A post-analysis implementation uses a guaranteeing fixed batch on the
4,371 eligible tasks and the adaptive rule elsewhere. It preserves precision
success at every completion while removing the need for feedback on
88.95\% of tasks. Both versions reach 98.64\% recorded success; mean width
changes from 18.59\% to 18.72\% of $W_0$. This saves a dependency between
answers, not readings, and is not a measured latency saving.

Exact rational checks reconstruct every batch and adaptive envelope.
Raw-source replay reproduces all 24,570 policy evaluations and 49,140 batch
intervals with no containment or success discrepancy. Independent endpoint
LPs and scale-aware tolerances are documented in S6; source calendars,
integer counter subtraction and all 6,552 frozen forecast rows are checked.
The geometric separation is proved on the stated record model; these
checks do not establish physical meter accuracy.
\paragraph{Sensitivity to additional hourly information.}
Before computing this sensitivity, we fixed the subset to all 290
original feedback-essential cases and the illustrative caps
$0\leq x_h\leq E/2$ and $0\leq x_h\leq E/3$.
These are scenario constraints, not validated physical limits. The primary comparison
keeps the original absolute requirement $\tau=W_0^{\rm old}/2$.
A separate calculation halves the new constrained width and therefore
asks a different precision question (Table~\ref{tab:counter-caps}).

\begin{table}[htbp]\centering\small
\caption{Illustrative-cap sensitivity of the fixed 290-case subset.
Each row partitions all cases into a fixed-batch guarantee, a
feedback-essential deterministic guarantee, and no deterministic
two-query guarantee. ``Old'' and ``new'' identify the initial width used
in the threshold, not different data.}
\label{tab:counter-caps}
\begin{tabular}{llrrr}\toprule
Hourly cap & Precision threshold & Batch & Feedback & Unresolved\\\midrule
None & $W_0^{\rm old}/2$ & 0 & 290 & 0\\
$E/2$ & $W_0^{\rm old}/2$ & 282 & 8 & 0\\
$E/3$ & $W_0^{\rm old}/2$ & 290 & 0 & 0\\
$E/2$ & $W_0^{\rm new}/2$ & 163 & 33 & 94\\
$E/3$ & $W_0^{\rm new}/2$ & 221 & 0 & 69\\\bottomrule
\end{tabular}\end{table}

At the unchanged requirement, the $E/2$ cap leaves eight feedback-essential
cases, five on meter~1 and three on meter~2; all retain exact randomized-batch
success $1/2$. The $E/3$ cap leaves none; 34 cases already satisfy the
old threshold without a query. The constructed example
$E=1$, $c=(0,0,1,1,0,1)$ illustrates the mechanism: imposing $x_h\leq1/2$
changes the best fixed-batch guarantee from one to $1/2$, matching
adaptation, while its initial width remains one.
Thus the original guarantee gap is sensitive to additional valid hourly
information; it is not established as robust to physical capacity bounds.

After freezing the capped geometry, an exact integer-register check found
that the recorded trace violates the $E/2$ assumption in 8 of 290 cases
(meters 1--3: 1, 0, 7) and the $E/3$ assumption in 57 cases (11, 8, 38).
These correspond to 8 and 48 of the 228 distinct blocks. All cases remain
in Table~\ref{tab:counter-caps}; the capped guarantees do not cover an
excluded trace. Compatibility does not establish a physical operating bound.
None of the eight remaining feedback-essential cases at the unchanged
threshold violates the $E/2$ cap in its recorded trace; the eight excluded
cases form a different set.

All results use the original coefficients and total. Exact rational
bounded-allocation formulas and independently assembled polytope LPs
agree across every case, batch and first query; the largest discrepancy
is $6.11\times10^{-16}$ in units of $W_0^{\rm old}$.
Finite adverse games use only exactly feasible capped profiles.
Their upper bounds are called exact continuum optima only when a
matching fixed-batch or two-batch guarantee is supplied.
The fixed protocol, 1,740-row ledger, coefficients and checkable witnesses
are included in Online Resource~1. This calculation does not estimate
full-panel prevalence or measured acquisition savings.

\FloatBarrier
% Optional concise insertion; the main manuscript is not edited by this agent.
\subsection{Replication on an external stored-counter source}
\label{sec:external-hpm}
We froze a public-source protocol for HeatpumpMonitor.org \citep{hpmabout,hpmapi} before fitting
forecasts or evaluating acquisition outcomes. The ascending-ID screen
required public MID status, explicit Class-1 electrical-meter metadata and
a cumulative electrical-energy feed spanning January--July 2025. The first
three qualifiers were systems 1, 2 and 7: systems 3 and 6 lacked MID status,
while 4 and 5 failed calendar coverage. The source rule was locked before
hourly retrieval. Detailed training gates were fixed after training-only
inspection but before July retrieval; all three systems met the requirements
of at least 95\% valid hourly intervals and 90 complete training days, without
replacement. Bundled metadata, request logs and protocol locks expose this
chronology and every exclusion. January--June 2025 supplied training data; July supplied 31 test
days per system. The four existing model specifications were fitted once,
without retuning or replacement. All three six-hour blocks and six model
pairs were retained, giving 1,674 tasks with no reference-data attrition
and no zero-width case. The half-initial-width target, five optional
internal cumulative samples and two-sample budget were unchanged.

The target here is the difference of consecutive hourly \emph{API-grid
stored cumulative samples}. The interface can label a selected native
sample with a grid timestamp \citep{hpmcode}. Native acquisition and historical delivery
times are unavailable, so chronological forecast features are established
on the declared grid, not as measured midnight production availability.
The exact stored-sample model does not assert exact physical energy or
unconditional Class-1 pointwise error bounds.

The exact triage found 1,614 tasks with a guaranteeing fixed batch, 29
requiring feedback for a deterministic half-width guarantee, and 31
without such a two-query adaptive guarantee. The 29 cases span all three
systems (2, 14 and 13 cases) and all six model pairs. Each has matching
rational adverse-profile and all-completion covering witnesses proving
optimal randomized-batch worst-case success exactly $1/2$; adaptation
guarantees success one. All 29 batch optima are exactly $1/2$, including
their minimum, median and quartiles.

\begin{table}[htbp]
\centering\small\setlength{\tabcolsep}{4pt}
\caption{External stored-sample replay. Success is the probability of
reaching half initial width. Randomized outcomes average the declared
mixture exactly over its finite support. The cover policy is specified
only on the 29 cases with certified batch optimum $1/2$. Full recovery
uses all five optional samples, outside the two-sample budget; its actual
API request cost is compared in Section~\ref{sec:actual-api}.}
\label{tab:external-hpm}
\begin{tabular}{lrrr}
\toprule
Policy & All-task success (\%) & 29-case success (\%) &
29-case mean $W/W_0$\\
\midrule
Exact adaptive & 99.223 & 100.000 & 0.3113\\
Best deterministic worst-width batch & 99.104 & 96.552 & 0.3736\\
Greedy adaptive & 99.104 & 100.000 & 0.3114\\
Mean-width-optimal batch & 95.250 & 82.004 & 0.4003\\
Precision-optimal covering batch & --- & 81.034 & 0.3942\\
Uniform batch & 63.053 & 34.138 & 0.5899\\
Full stored-trace recovery & 100.000 & 100.000 & 0.0000\\
\bottomrule
\end{tabular}
\end{table}

In recorded replay, exact adaptation and the best deterministic batch
both succeed in 1,658 tasks; only adaptation succeeds in three, only the
batch in one, and neither in twelve. The net observed gain is therefore
two tasks, or 0.119 percentage points. Ordering certification is 91.219\%
versus 90.741\%, with 14 adaptive-only and six batch-only decisions.
The replication extends the information-model separation to a new source;
it does not demonstrate a large improvement over an ordinary adaptive
policy or the provider's permitted full-range retrieval. These dependent
tasks represent three systems, not 1,674 independent deployments.
Section~\ref{sec:actual-api} compares that full-range alternative under
the actual provider interface and measured request times.

A separate internal checker verifies 16,740 terminal intervals by exact
rational arithmetic, all 29 matching half-success witnesses, and the
two-side adaptive formula for every positive-width task. Independent
fiber LPs for the earliest task in each system--pair group (360 solves)
agree with the interval endpoints to $2.78\times10^{-17}$; solver
feasibility tolerance is $10^{-9}$ and comparison tolerance $2\times10^{-8}$.
All 2,232 saved persistence values equal the corresponding prior-day
grid differences. This is independent implementation checking within
the study, not an external human validation.

\FloatBarrier

\paragraph{Measurement uncertainty and a sufficient reading budget.}\label{sec:stored-error-sensitivity}
We retain all 29 external half-separation cases and ask whether feedback's
assurance advantage survives errors in the retrieved internal readings.
Normalize the retained total to one and divide the original coefficient
vector by its positive range, without translation. Certificate-LP widths
are then measured in units of the original $W_0$, with initial width one.
The hidden record consists of
$y_h\geq0$, $\sum_h y_h=1$, and five fixed errors
$e_k\in[-\delta,\delta]$; query $k$ returns
$\sum_{h\leq k}y_h+e_k$. Errors are fixed before policy randomization.
The endpoints remain exact, and no hourly cap is imposed in this separate
sensitivity calculation. Before solving, we fixed
$\delta\in\{0,0.001,0.01\}$ and the unchanged requirement $W\leq W_0/2$.
Initial width remains $W_0$, so halving the new initial width gives the
same threshold. These allowances concern block-total units; they are
illustrative bounds, not calibrated physical-device specifications.

\begin{table}[htbp]\centering\small
\caption{Guarantees certified on the fixed 29-case external subset under
bounded internal-readout errors and exact endpoints. Every passing case
has adaptive worst-case success one and optimal randomized-batch worst-case
success exactly $1/2$. All six model pairs are represented at each level.
Exact lower witnesses rule out a deterministic two-query guarantee in the
remaining cases.}
\label{tab:stored-error-certificates}
\begin{tabular}{lrrrr}\toprule
Error allowance & Certified cases & System 1 & System 2 & System 7\\\midrule
$0$ & 29/29 & 2/2 & 14/14 & 13/13\\
$0.001E$ & 27/29 & 2/2 & 13/14 & 12/13\\
$0.01E$ & 22/29 & 2/2 & 12/14 & 8/13\\\bottomrule
\end{tabular}\end{table}

Two solver-free checkers verify all 49 positive-error certificates using
exact rational weak duality and an explicit correction for residuals in
the proposed multipliers. The smallest certified margin below $W_0/2$
exceeds $0.003759W_0$. The complete LPs, multipliers, finite variable bounds
and derivation are in Online Resource~1, S6. These are upper guarantees
for the saved first-query/successor supports, not claims of exact
optimal-width values.

The original batch upper bound $1/2$ survives because zero-error records
and fibers remain feasible. The two-successor guarantee supplies the
matching continuum lower bound. Thus the exact separation survives in
27 and 22 cases, spanning every system and model pair. Independent
error-eliminated LPs reproduce all 87 case--level calculations with maximum
discrepancy $1.89\times10^{-15}$.

For the remaining two cases at $0.001E$ and seven at $0.01E$, every legal first query $k$ has a
rational answer $t_k$ such that, for each possible second query $j$, two
feasible latent records share first answer $t_k$ and an identical second
answer, while their target contrast exceeds $W_0/2$. Every deterministic
two-query policy therefore has a permitted failing record; stopping early
cannot help. Two exact checkers verify 45 first-answer witnesses and 180
endpoint pairs, with minimum strict excess $5.9060\times10^{-6}W_0$.
The loss of the guarantee is caused by larger observational fibers, not
LP tie breaking. No randomized-adaptive impossibility is asserted.

\paragraph{A certified budget fallback covering all 29 cases.}
A post-analysis calculation keeps the cases, $0.01E$ allowance, exact
endpoints and precision threshold, and permits a third reading. Among
899 fixed-batch LPs, exact upper certificates identify a three-reading
batch for every case; the largest upper bound is $0.458896W_0$.
The implemented rule chooses two adaptive readings for 22 cases and a
fixed three-reading batch for seven, before inspecting optional answers.
It guarantees precision in all 29 cases.

Its two-query branches minimize certified affine upper bounds by exact
rational comparisons, without online LPs. Their lower envelope is at most
$W_0/2$ on the whole first-answer interval $[-0.01,1.01]$.
Separate tests cover 1,015 constructed executions and reject incompatible
transcripts. The continuum guarantee comes from the certificates, not
these tests. This executable branch rule can differ from the earlier
conditional-width minimizer, so its recorded replay figures are not reused.

\begin{table}[htbp]\centering\small
\caption{Post-analysis budget fallback on the unchanged 29-case external
subset at $0.01E$ internal error. Totals sum hard reading allowances across
separate pairwise audits; they are not measured API requests or costs.}
\label{tab:stored-error-fallback}
\begin{tabular}{lrrr}\toprule
Policy & Guaranteed cases & Maximum budget & Total allowances\\\midrule
Two-query adaptive & 22/29 & 2 & 58\\
Minimum fixed-batch budgets & 29/29 & 3 & 87\\
Certified two-or-three rule & 29/29 & 3 & 65\\\bottomrule
\end{tabular}\end{table}

These are the minimum per-case deterministic hard budgets: a one-query
policy is a fixed batch and cannot guarantee the threshold on any of these
cases; the seven lower witnesses additionally rule out two-query policies.
The upper certificates attain the corresponding budgets of two and three.
Every nonadaptive batch with at most two queries still has a failing
completion, so randomization cannot give sure success using any such batch
with positive probability. Thus three readings per case also attain the
minimum nonadaptive budget for sure success. The reduction from 87 to 65 is
$22/87=25.29\%$ in reserved readings on this selected subset. It does not
optimize expected stopping costs, share readings across model pairs, or
change the original two-query result; full-range retrieval remains the
preferred option on the measured public API.
\FloatBarrier
\subsection{Retrieval interface and measured operational consequences}
\label{sec:actual-api}
The public source permits both single-time and full-range queries
\citep{hpmapi,hpmcode}. We therefore tested the latter as an operational
comparator instead of treating a two-reading budget as a measured tariff.
The first eligible daily-persistence--MLP task from each external system
was fixed before retrieving optional answers. Two repetitions of five
policies produced 30 trials and 60 actual HTTP requests, including six
common endpoint reads. The protocol limited request starts to one per
second, with no retries or cache-busting. All trials succeeded; timings
include the recorded persistent-connection behavior and uncontrolled
public cache/server state.

\begin{table}[htbp]\centering\small
\caption{Actual public-API measurements on three fixed tasks, two
repetitions per policy. Incremental I/O excludes the separately logged
client throttle wait. Bytes are response-body bytes, not wire traffic.
The full-range interface returns seven hourly samples in one request,
including retained endpoints; it is an operational alternative outside
the two-optional-sample benchmark budget.}\label{tab:actual-api}
\begin{tabular}{lrrr}\toprule
Policy & Requests & Mean I/O (s) & Mean body bytes\\\midrule
Exact adaptive & 2 & 0.497 & 104.3\\
Best deterministic worst-width batch & 2 & 0.539 & 104.7\\
Mean-width-optimal batch & 2 & 0.582 & 104.3\\
Greedy adaptive & 2 & 0.551 & 104.3\\
Full hourly range & 1 & 0.331 & 226.7\\\bottomrule
\end{tabular}\end{table}

All five policies attained the precision target on these three traces.
Bulk retrieval used fewer requests and less observed I/O time, while
also reducing terminal width to zero. Selective retrieval reduced body
size by about 122 bytes; this small difference excludes HTTP/TLS overhead
and does not establish an operational saving. A range returning only the
five internal samples could reduce the bulk response further. Including
the imposed client throttle, mean completion times were 1.984 seconds
for exact adaptation and 1.040 seconds for bulk retrieval; these are not
intrinsic provider latencies. The small workload does not support a
population speedup or monetary-cost claim.

An offline checker matches all 60 received payload hashes to separately
cached source responses, checks all 12 adaptive second-query decisions
and all 30 interval containments (maximum discrepancy
$1.74\times10^{-18}$). Each result can be verified without repeating
provider traffic. The appropriate decision for this observed interface
is bulk retrieval when its unrestricted response is admissible. The
two-query separation instead answers the distinct question posed by a
fixed sample budget; it does not override an available cheaper interface.
\FloatBarrier
\paragraph{Selecting the acquisition protocol.}
The interface, budget unit and operational preference must be specified before
an acquisition guarantee is used. A reference implementation selects among
the supplied fixed, adaptive and full-range certificates for the 29 cases.
It checks permitted
queries, scalar-reading and information-round caps, transport grouping, and
the requested preference; a missing guarantee triggers escalation rather than
a relaxed target. This is an integration of the preceding results, not a new
optimization theorem. A frozen offline evaluation retains all 29 external
cases, the exact endpoint total, the $0.01E$ stored-readout allowance and the
unchanged half-width target.

\begin{table}[t]
\centering\small
\caption{Protocol selection under four explicitly declared interfaces. The
restricted interfaces are benchmark assumptions; the public API permits bulk
retrieval. All rows use the same 29 cases and uncertainty model.}
\label{tab:protocol-selection}
\begin{tabular}{p{0.40\linewidth}cp{0.36\linewidth}}
\hline
Interface and preference & Certified & Selected protocol \\
\hline
Range available; fewest requests & 29/29 & Bulk range in all 29 \\
Three readings, two rounds; fewest readings & 29/29 & Adaptive two in 22; fixed three in seven \\
Two readings, two rounds & 22/29 & Adaptive two in 22; escalate seven \\
Three readings, one round & 29/29 & Fixed three in all 29 \\
\hline
\end{tabular}
\end{table}

Unrestricted range retrieval is preferred on the actual public interface.
Its five internal values contain each case's certified three-query subset,
so information monotonicity preserves the precision guarantee under bounded
stored-readout errors; it does not imply zero uncertainty about true increments.
The restricted three-reading rule reserves 65 scalar observations in 51
information rounds, whereas the one-round-per-case rule reserves 87 in 29
rounds. These are distinct resources. HTTP counts are left unspecified for
the hypothetical restricted interfaces. In particular, an interface offering
a single native multipoint request can favor fixed-three under request
priority while favoring adaptive-two under reading priority.

The preceding three-task API timings support bulk selection on that service;
they are not measurements on this 29-case sensitivity subset. The supplement exports all
116 case--interface decisions, exclusions and certificate hashes, with an
offline callback dispatcher and separate verification. The implementation returns an auditable acquisition plan: a supported
case, its measurement assumptions, selected query protocol, hard resource
allowances and the certificate proving its precision. If no supplied plan
meets the interface, it returns escalation. This catalogue decision alone
does not prove that every possible policy fails; the seven deterministic
two-query exclusions additionally have the exact lower witnesses described
above.

The operational wrapper binds each supplied certificate to its saved model
pair, exact coefficients, meter, API-grid block, retained endpoints and source
hashes; readout envelopes with a different task, timestamp, feed or unit are
rejected. For example, the system~1 ridge--MLP comparison on 6 July 2025,
12:00--18:00 UTC, uses boundary~2 followed by boundary~1 or~4 under the
restricted two-reading rule, while its public-interface plan selects bulk
retrieval. These are checks against accidental task mismatch and a worked
planning example; measured API trials and constructed dispatch tests remain
separate, and the wrapper does not authenticate a provider or calibrate a meter.
\paragraph{A reading--feedback cost frontier.}
A post-analysis calculation makes the protocol choice explicit without
assigning an unmeasured monetary benefit. Let $c_{\rm read}>0$ charge for a
reserved scalar reading and $c_{\rm round}\geq0$ for an information round;
write $r=c_{\rm round}/c_{\rm read}$. Under the restricted interface, each of
the 22 cases admitting the certified adaptive rule costs $2+2r$ in reading
units, versus $3+r$ for its certified fixed batch. Thus feedback is preferable
exactly when $r<1$. The seven remaining cases use fixed-three under both
complete-coverage baselines. Over the same 29 cases, fixed-three costs
$87+29r$, adaptive-where-certified with fixed fallback costs $65+51r$, and
the catalogue selector attains their minimum. Ties favor fewer rounds.

\begin{table}[htbp]\centering\small
\caption{Declared-cost sensitivity on the unchanged 29-case subset at 1\%
internal error and exact endpoints. Costs are reading-charge units, not
measured prices or latency. All methods retain the same precision guarantee.}
\label{tab:protocol-cost-frontier}
\begin{tabular}{rrrr}\hline
$r$ & Fixed-three & Adaptive with fixed fallback & Catalogue choice\\\hline
0 & 87 & 65 & 65\\
0.25 & 94.25 & 77.75 & 77.75\\
0.5 & 101.5 & 90.5 & 90.5\\
1 & 116 & 116 & 116\\
2 & 145 & 167 & 145\\
4 & 203 & 269 & 203\\\hline
\end{tabular}\end{table}

Neither baseline is uniformly inferior: fixed-three is already optimal
within this catalogue when $r\geq1$, and adaptive with fallback is optimal
when $r\leq1$. The benefit is preventing an acquisition choice that assigns
priority to the wrong resource. The analysis uses exact rational arithmetic,
exports all 174 case--ratio choices with their certificates, and independently
checks all 306 candidate evaluations. It is separable across pairwise audits;
shared readings, stopping costs and unlisted protocols are outside this
calculation. The existing three-task public-API measurements remain separate:
mean I/O was 0.331~s for bulk and 0.497~s for adaptive retrieval, supporting
bulk on that admissible interface. Those observations do not calibrate $r$
or turn the modeled frontier into measured savings.
\subsection{Computability and proof checks}
The proofs establish the universal classification. Numerical checks address
the implementation separately: 3,202 endpoint LPs and 28 finite games agree
with the analytic constructions, with largest width discrepancy
$3.34\times10^{-16}$ against a $2\times10^{-8}$ scale-aware tolerance.
Six explicit polytopes check overlapping constraints, signed contrasts,
boxes and coupled capacities. All 37 enumerated batches and 24 possible
first queries agree with the corresponding LP guarantees. These are small
formulation checks, not a general scalability claim.

For coordinate queries under disjoint exact totals and unit costs, a
supporting compact mean-width solver avoids enumerating batches.
Its retained-endpoint flow representation and feasible dual completion
and adaptive lower-bound proof are given in Online Resource~1, S5;
S7 supplies the measured comparisons. General weighted-oracle
and dynamic-program LP constructions are established
\citep{kawase2024,martin1990}. The specialized representation and its
hard-budget policy decomposition are the implementation contribution.

All 28 archived compact settings solved. Twelve prescribed enumerated
comparisons agreed within $1.27\times10^{-15}$ after normalization.
Compact was faster in eight and used less whole-process peak memory in
six; enumeration was faster in the four smallest one-block settings.
The largest tested LP had 38,820 variables and ran in 9.933--10.012 seconds
for 48 target coordinates. These are single-run measurements on a shared
Intel Core i7-8565U host with one numerical thread. The complete timing,
memory, residual, tie-profile and skip ledgers are in S1/S7.
No large-scale optimizer for arbitrary convex sets or precision risk is implied.
\subsection{Historical controls: guarantees and observed accuracy}
\label{sec:historical-application}
The earlier completed-archive study \citep{portugal,uci} incorporates every permitted forecast
copy before evaluation. Under identical information and zero new readings,
paired bounds certify 7,836/9,720 Portugal comparisons (80.62\%), versus
2,513 (25.85\%) from separate-loss bounds. However, a donor imputer makes
no observed whole-period errors at full coverage; the added value is an
all-completions certificate, not demonstrated prediction accuracy.
In the daily archive analysis, free reconstruction already resolves 95.33\%
of masked targets, exact acquisition adds only 0.042 percentage points over
greedy, and robust randomization worsens realized width despite improving
its objective. Outcome-independent optimal-face refinements retain that
tradeoff \citep{iancu2014}. The earlier 0.9738064\% mean acquisition saving
and its denominator remain unchanged. S2--S4 retain these adverse controls,
all 324 grid cells, four-model gates, bootstrap details and full protocols.
They explain why the present paper evaluates the assurance requirement and
retrieval interface explicitly rather than claiming general acquisition savings.
\section{Discussion}\subsection{What the method changes in an audit}
The practical output is a justified acquisition choice. The auditor first
states the contrast to compare, the precision required, the information
already available, and the retrieval options. A sufficient initial interval
needs no query. Otherwise, a bulk response or fixed batch can avoid feedback;
adaptive retrieval is useful when it meets the required guarantee with fewer
permitted readings. This ordering prevents a sophisticated selective policy
from being used where a simpler service already supplies adequate information.

The public API is a concrete example. Its full-range response is admissible,
uses one request, and is faster on the measured three-task workload. The
selector therefore recommends bulk. Under the separate restricted-reading
contract, the same assurance analysis instead justifies 22 two-reading and
seven three-reading prescriptions. The contracts limit different resources:
65 versus 87 counts reserved scalar readings across the selected 29 audits,
while 51 versus 29 counts information rounds. A service that values rounds
more than readings can reasonably choose fixed batches. Actual monetary
costs and service-wide latency require corresponding operational measurements.

\paragraph{Breadth of the theory and limits of the implementations.}
The classification applies to any known compact convex feasible set, scalar
affine target, finite exact-linear query dictionary and answer-independent
family of legal final query sets. It permits signed targets, overlapping
constraints, valid upper bounds and unequal fixed costs. It is not restricted
to energy or disjoint totals. Fixed bounded errors also fit by including
them in the hidden record, as detailed in Online Resource~1, S6.
The finite dictionary matters: allowing every linear functional at equal
cost would permit a direct query of the scalar contrast itself.

Computational guarantees have narrower scopes. The general polyhedral
formulations handle a declared batch family and exact deterministic
two-query planning. The counter rule exploits cumulative sums and a known
nonnegative total. The compact coordinate solver uses disjoint exact totals
and unit costs. These efficient specializations do not supply an efficient
arbitrary-budget precision optimizer for every convex set.

\paragraph{What the application establishes.}
Two public sources exhibit the same exact feedback-dependent assurance gap.
Its prevalence is 5.90\% and 1.73\%, and overall recorded gains over the best
fixed deterministic worst-width batch are only 0.407 and 0.119 percentage
points. Greedy adaptation succeeds on all 29 external separation cases.
Thus the main benefit is identifying when a guarantee needs feedback, not
a large observed improvement over every adaptive heuristic. Three systems
per source, repeated model pairs and six-hour blocks do not establish
population prevalence.

Illustrative capacity constraints often make feedback unnecessary, while
measurement errors can increase the sufficient reading budget. The
sensitivity uses exact endpoints and declared internal errors; deployment
requires defensible error and physical bounds. The half-initial-width target
is a declared precision requirement, not a calibrated financial threshold.
A narrow interval can still contain zero. Model ordering, shared retrievals
across comparisons, and requests until a decision are separate objectives.

\section{Conclusions}
The risk objective determines whether feedback can improve the worst-case
guarantee for a model-comparison interval. Concavity gives a complete
no-improvement classification; nonconcave penalties admit sharp query-based
separations. An exact two-query rule and checkable certificates turn that
distinction into an acquisition decision. The reference implementation chooses among supplied bulk, fixed and adaptive
certificates using declared access and resource limits. New comparisons require
new certificates; a catalogue exclusion alone is not an impossibility proof.
The public-data results show both when feedback changes assurance and when
bulk or fixed retrieval is preferable. The practical contribution is this
auditable choice, supported by explicit assumptions and executable evidence.
\section*{Supplementary Materials}
Online Resource 1, \texttt{research\_supplement.zip}, contains the editable
manuscript, executable methods, frozen forecasts, source identifiers, case
ledgers, exact witnesses and numerical checks. Start with \texttt{README.txt};
\texttt{CURRENT\_RESULTS\_MAP.txt} maps each current table and verification
claim to its command, inputs and outputs. The default verification replays
saved forecasts and does not train new models. Separate commands cover
raw-input regeneration and fresh fitting. BLEM-derived traces are reconstructed
locally; attributed HeatpumpMonitor records are bundled. The indexed extended
methods retain full proofs, solver scaling, historical controls, forecasting
protocols and uncertainty checks. Protocol-selection files include all 116
interface decisions and task-bound readout checks. The separate checker
compares saved contrasts with frozen geometry, verifies forecast hashes,
and reconstructs retained endpoints from cached API records. Historical
reports remain available as clearly distinguished development records.

\section*{Additional Information and Declarations}
\subsection*{Data and Code Availability}
The numerical code and detailed outputs are supplied in Online Resource 1.
The public source identifiers are given in the references. The package
separates bundled Portugal/household replay from BLEM raw-source
regeneration; it does not redistribute the latter's raw counter traces.
Constructed polytopes include their complete inequalities and objectives.
Dual outputs retain raw multipliers, any numerical mass correction, the
completed record, and pricing and enumerated lower bounds. LP and endpoint
cross-checks use floating-point arithmetic with explicit tolerances; the
counter covering and adverse-profile certificates are additionally checked
with exact rational arithmetic on their serialized inputs. No public
repository or persistent archived release is claimed.

\subsection*{Funding}
This research received no external funding.

\subsection*{Competing Interests}
The author declares no competing interests.

\subsection*{Author Contributions}
Jishan Kapoor defined the research objectives and evaluation requirements,
directed and supervised the analysis and revisions, interpreted the findings,
and approved the manuscript. The computational and drafting assistance used
under this supervision is described below.

\subsection*{Research Transparency}
This work uses public data and constructed mathematical instances, with
no newly recruited participants. New measurements concern public API
response times and payloads, with their small-workload scope reported. Retrospective masks and
archive-access assumptions are identified explicitly. Comparisons use
the same information and the same stated hard budgets. The new proofs
and numerical checks do not change the historical forecast results.

\subsection*{Generative-AI Assistance}
Jishan Kapoor provided the topic, set the research objectives and evaluation
requirements, directed the analyses and tests, and closely supervised all
use of OpenAI Codex (OpenAI; accessed October 2026; model versions were not
consistently recorded). The tool assisted with literature searches,
development and checking of research content and code, computational
execution, and manuscript drafting and editing. It supported iterative
verification and clearer presentation. The author reviewed and revised
these outputs and retains full responsibility for the research decisions,
interpretation of results, and final manuscript. Separate checker
implementations within this study are not external independent replication.
This generative assistance is distinct from the boosted-tree and neural
forecasting models fitted to the public application data.


\FloatBarrier
\begin{thebibliography}{99}
\bibitem[Amenta et~al.(2016)]{amenta2016}
Amenta N, De Loera JA, Sober\'on P (2016) Helly's theorem: new variations
and applications. Author manuscript, arXiv:1508.07606v2.
\url{https://arxiv.org/abs/1508.07606v2}

\bibitem[Gaillard et~al.(2023)]{gaillard2023}
Gaillard, P., Gerchinovitz, S., \& de Montbrun, {\'E}. (2023).
Adaptive approximation of monotone functions [Preprint].
arXiv. \url{https://arxiv.org/abs/2309.07530v1}

\bibitem[Hara et~al.(2024)]{hara}
Hara, S., Matsuura, M., Honda, J., \& Ito, S. (2024).
Active model selection: A variance minimization approach.
\emph{Machine Learning, 113}, 8327--8345.
\url{https://doi.org/10.1007/s10994-024-06603-1}

\bibitem[Hebrail and Berard(2006)]{uci}
Hebrail, G., \& Berard, A. (2006). Individual household electric power
consumption [Data set]. UCI Machine Learning Repository.
\url{https://doi.org/10.24432/C58K54}

\bibitem[Iancu and Trichakis(2014)]{iancu2014}
Iancu, D. A., \& Trichakis, N. (2014). Pareto efficiency in robust optimization.
\emph{Management Science, 60}(1), 130--147.
\url{https://doi.org/10.1287/mnsc.2013.1753}

\bibitem[Kawase and Sumita(2024)]{kawase2024}
Kawase, Y., \& Sumita, H. (2024). Randomized strategies for robust
combinatorial optimization with approximate separation.
\emph{Algorithmica, 86}, 566--584.
\url{https://doi.org/10.1007/s00453-023-01175-3}

\bibitem[Kossen et~al.(2021)]{active}
Kossen, J., Farquhar, S., Gal, Y., \& Rainforth, T. (2021).
Active testing: Sample-efficient model evaluation.
\emph{Proceedings of Machine Learning Research, 139}, 5753--5763.
\url{https://proceedings.mlr.press/v139/kossen21a.html}

\bibitem[Kostmann and H\"ardle(2019)]{blem}
Kostmann, M., \& H\"ardle, W. K. (2019). Forecasting in blockchain-based
local energy markets. \emph{Energies, 12}, 2718.
\url{https://doi.org/10.3390/en12142718}
Source data: \url{https://github.com/QuantLet/BLEM}.

\bibitem[Krieg et~al.(2025)]{krieg2025}
Krieg, D., Novak, E., \& Ullrich, M. (2025).
On the power of adaption and randomization.
\emph{Forum of Mathematics, Sigma, 13}, e152.
\url{https://doi.org/10.1017/fms.2025.10101}

\bibitem[Krieg and Ullrich(2026)]{krieg2026}
Krieg, D., \& Ullrich, M. (2026).
Approximation of functions: Optimal sampling and complexity.
\emph{Acta Numerica, 35}, 273--457.
\url{https://doi.org/10.1017/S0962492925100287}
Author manuscript: \url{https://arxiv.org/abs/2602.02066v2};
section and equation locators in the text refer to that version.

\bibitem[Kuboyama(2026)]{kuboyama2026}
Kuboyama, T. (2026).
How many labels does model choice need? Certificates and budgets for selective
prediction [Preprint].
arXiv. \url{https://arxiv.org/abs/2609.18622v1}

\bibitem[Kurtz(2026)]{kurtz2026}
Kurtz, J. (2026). Bounding the optimal number of policies for robust
K-Adaptability. \emph{Mathematical Programming}. Published 29 January 2026.
\url{https://doi.org/10.1007/s10107-026-02329-1}

\bibitem[Maia Polo et~al.(2024)]{polo}
Maia Polo, F., Maity, S., Yurochkin, M., Banerjee, M., \& Sun, Y. (2024).
Weak supervision performance evaluation via partial identification.
\emph{Advances in Neural Information Processing Systems, 37}.
\url{https://doi.org/10.52202/079017-4309}

\bibitem[Manski(2005)]{manski}
Manski, C. F. (2005). Partial identification with missing data:
Concepts and findings. \emph{International Journal of Approximate
Reasoning, 39}, 151--165.
\url{https://doi.org/10.1016/j.ijar.2004.10.006}

\bibitem[Martin et~al.(1990)]{martin1990}
Martin, R. K., Rardin, R. L., \& Campbell, B. A. (1990).
Polyhedral characterization of discrete dynamic programming.
\emph{Operations Research, 38}(1), 127--138.
\url{https://doi.org/10.1287/opre.38.1.127}

\bibitem[Novak and Wo{\'z}niakowski(2008)]{novak2008}
Novak, E., \& Wo{\'z}niakowski, H. (2008).
\emph{Tractability of multivariate problems. Volume I: Linear information}.
European Mathematical Society.
\url{https://doi.org/10.4171/026}

\bibitem[Novak(1995)]{novak1995}
Novak E (1995) The adaption problem for nonsymmetric convex sets.
Journal of Approximation Theory 82(1):123--134.
\url{https://doi.org/10.1006/jath.1995.1071}

\bibitem[Okanovic et~al.(2025)]{okanovic2025}
Okanovic, P., Kirsch, A., Kasper, J., Hoefler, T., Krause, A., \&
G\"urel, N. M. (2025). All models are wrong, some are useful:
Model selection with limited labels.
\emph{Proceedings of Machine Learning Research, 258}, 2035--2043.
\url{https://proceedings.mlr.press/v258/okanovic25a.html}

\bibitem[OpenEnergyMonitor(2026a)]{hpmabout}
OpenEnergyMonitor (2026a) HeatpumpMonitor.org: About.
\url{https://heatpumpmonitor.org/about}. Accessed 12 October 2026.

\bibitem[OpenEnergyMonitor(2026b)]{hpmapi}
OpenEnergyMonitor (2026b) HeatpumpMonitor.org API helper.
\url{https://heatpumpmonitor.org/api-helper}. Accessed 12 October 2026.

\bibitem[OpenEnergyMonitor(2026c)]{hpmcode}
OpenEnergyMonitor (2026c) HeatpumpMonitor.org timeseries controller and
Emoncms PHPTimeSeries implementation. Source snapshots and hashes supplied
in Online Resource 1.
\url{https://github.com/openenergymonitor/heatpumpmonitor.org/blob/main/www/Modules/timeseries/timeseries_controller.php};
\url{https://github.com/openenergymonitor/emoncms/blob/master/Modules/feed/engine/PHPTimeSeries.php}.
Accessed 12 October 2026.

\bibitem[Osipenko(2019)]{osipenko}
Osipenko, K. Yu. (2019). Generalized adaptive versus nonadaptive recovery
from noisy information. \emph{Journal of Complexity, 53}, 162--172.
\url{https://doi.org/10.1016/j.jco.2018.12.001}

\bibitem[Patton(2011)]{patton}
Patton, A. J. (2011). Volatility forecast comparison using imperfect
volatility proxies. \emph{Journal of Econometrics, 160}(1), 246--256.
\url{https://doi.org/10.1016/j.jeconom.2010.03.034}

\bibitem[Rockafellar(1965)]{rockafellar1965}
Rockafellar RT (1965) Helly's theorem and minima of convex functions.
Duke Mathematical Journal 32:381--397.
\url{https://sites.math.washington.edu/~rtr/papers/rtr005-Helly.pdf}

\bibitem[Sawade et~al.(2012)]{sawade}
Sawade, C., Landwehr, N., \& Scheffer, T. (2012).
Active comparison of prediction models.
\emph{Advances in Neural Information Processing Systems, 25}.
\url{https://proceedings.neurips.cc/paper_files/paper/2012/file/92fb0c6d1758261f10d052e6e2c1123c-Paper.pdf}

\bibitem[Sion(1958)]{sion1958}
Sion, M. (1958). On general minimax theorems.
\emph{Pacific Journal of Mathematics, 8}(1), 171--176.
\url{https://doi.org/10.2140/pjm.1958.8.171}

\bibitem[Trindade(2015)]{portugal}
Trindade, A. (2015). ElectricityLoadDiagrams20112014 [Data set].
UCI Machine Learning Repository.
\url{https://doi.org/10.24432/C58C86}

\bibitem[Vayanos et~al.(2026)]{ddid}
Vayanos, P., Georghiou, A., \& Yu, H. (2026).
Robust optimization with decision-dependent information discovery.
\emph{Management Science, 72}(2), 1509--1528.
\url{https://doi.org/10.1287/mnsc.2021.00160}

\bibitem[Wei and Zhang(2024)]{wei2024}
Wei, N., \& Zhang, P. (2024). Adjustability in robust linear optimization.
\emph{Mathematical Programming, 208}, 581--628.
\url{https://doi.org/10.1007/s10107-023-02049-w}
% New entries; place Gaillard before Hara and Kuboyama after Krieg in the
% alphabetized bibliography. Existing wei2024 and ddid entries are reused.
\end{thebibliography}
\end{document}
